\section{Introduction}
\label{sec:introduction}

%\todo[inline]{R2 on queries for outliers: The stream
%database is assumed to represent
%deviating/abnormal/outlier behavior,
%so that the discovered queries capture the common
%patterns that are inherent to the respective
%situations.
%We intend to clarify this in the paper.}
Complex event processing (CEP) is a computational paradigm
to facilitate reactive and proactive
applications~\cite{DBLP:journals/csur/CugolaM12,DBLP:journals/vldb/GiatrakosAADG20}.
 In CEP, a set of queries are evaluated continuously over
a stream of events in order to detect situations of
interest in domains, such as urban
transportation~\cite{DBLP:conf/edbt/ArtikisWSBLPBMKMGMGK14},
computational
finance~\cite{DBLP:journals/computer/ChandramouliAGSR10},
and supply chain
management~\cite{DBLP:journals/cii/KonovalenkoL19}.
Often, these situations of interest are deviating or
abnormal behaviour, so that their timely detection
facilitates the implementation of counter-measures.

In essence, CEP queries capture sequential patterns of
events, i.e., they describe a sequence of events of
specific types, additional conditions on the attribute
values of the respective events, and a time window for
their occurrence. As the exact characterization of the
event patterns is challenging in practice, one may support
users in the specification of queries
with techniques for automated query
discovery~\cite{icep,ilminer,DBLP:conf/btw/Kleest-Meissner23}.
 Here, the assumption is
that a database of finite, historic (sub-)streams that
cover the situation of interest is available. That is,
if the situation of interest is some deviating or
abnormal behaviour, the assumption is that the
database contains only the respective streams, not the
common behaviour. Then,
queries that match a certain share of these streams, as
determined by a support threshold, can be extracted and
subsequently be assessed and validated by a user.

\begin{figure}[t]
	\centering
	\includegraphics[width=0.8\columnwidth]{img/overview.drawio.pdf}
	%\vspace{-1.0em}
	\caption{The general setting of incremental event
	query discovery: An update of the stream
	database or support threshold shall be incorporated in
	an updated set of queries.}
	\label{fig:overview}
	\vspace{-1em}
\end{figure}

However, over time, the database of streams used as a
basis for query discovery may be subject to change.
Additional streams that capture a situation of interest
may become available, or streams that are already in the
database may be considered outdated. As part of that, the
requirements for query discovery in terms of the adopted
support threshold may also change, i.e., it may be relaxed
to increase the number of discovered queries, or
strengthened to avoid false positives in the result.

Yet, upon changes in the database or in the support
threshold, existing discovery algorithms need to construct
the resulting queries from scratch. This is problematic
since solving the problem of event query discovery is
computationally hard, as an exponentially growing space of
candidate queries (in the query length and attribute
domains) needs to be explored and the evaluation of a
single query also shows an exponential runtime complexity
(in the number of variables, i.e., conditions over
multiple events)~\cite{icdt2022}.

In this paper, we set out to avoid
inefficient recomputation with strategies for incremental
event query discovery. Our idea is that the discovery
result shall be reused and adapted once the databases or
the support threshold change, as illustrated in
\autoref{fig:overview}. To this end, after giving
formal preliminaries (\autoref{sec:problem-setting}), we
make the following contributions:
\begin{enumerate}[left=.2em,nosep,label=(\roman*)]
\item We present a theoretical analysis of
the problem context (\autoref{sec:updates}).
Specifically, we show that under certain conditions, the
result set of queries for an updated stream database or an
updated support threshold denotes a subset of the existing
result.
\item We present algorithmic
solutions for incremental query discovery
(\autoref{sec:algos}). They adapt a given set of queries,
thereby aiming to avoid inefficient recomputation.
\end{enumerate}
We assess the feasibility of our
solution in several experiments
(\autoref{sec:evaluation}), showing that incremental
discovery reduces the runtimes by up to 10x.
Finally, we review related work
(\autoref{sec:related-work}) and conclude the paper
(\autoref{sec:conclusions}).







